\documentclass[10pt,letterpaper]{article}

\usepackage[margin=0.5in]{geometry}
\usepackage[T1]{fontenc}
\usepackage[utf8]{inputenc}
\usepackage{lmodern}
\usepackage[hidelinks]{hyperref}
\usepackage{enumitem}
\usepackage{titlesec}

\pagestyle{empty}
\setlength{\parindent}{0pt}
\setlength{\parskip}{0pt}
\setlist[itemize]{leftmargin=12pt, itemsep=1pt, topsep=1pt, parsep=0pt, partopsep=0pt}

\titleformat{\section}{\large\bfseries}{}{0pt}{}[\titlerule]
\titlespacing*{\section}{0pt}{5pt}{2pt}

\newcommand{\resumeSubheading}[4]{%
  \textbf{#1} \hfill #2 \\
  #3 \hfill #4 \\
}

\begin{document}

\begin{center}
    {\Large \textbf{Sanskar Jain}} \\
    DevSecOps Engineer \;|\; CI/CD \;|\; Kubernetes \;|\; Terraform \;|\; AI Agents / MCP \\
    Hyderabad, India \enspace | \enspace 9021620156 \enspace | \enspace
    \href{mailto:sanskarjain0901@gmail.com}{sanskarjain0901@gmail.com} \enspace | \enspace
    \href{https://www.linkedin.com/in/sanskar-jain0901/}{LinkedIn}
\end{center}

\vspace{2pt}

\section*{Summary}
Software engineer with 2+ years delivering secure cloud services, containerised workloads, and automated \textbf{CI/CD} pipelines. Hands-on with \textbf{Azure DevOps}, Git, AWS, Kubernetes/EKS, Docker, Terraform, C\#, Python, and DevSecOps practices grounded in OAuth 2.0, JWT validation, and production observability. Built internal \textbf{AI agents} using the OpenAI API and \textbf{Model Context Protocol (MCP)} integrations.

\section*{Skills}
\textbf{DevSecOps \& Security:} Secure SDLC, OAuth 2.0, JWT, Okta, Kong Gateway, Input Validation, Secrets Management, Security Gates \\
\textbf{Cloud \& Containers:} AWS (EKS, ECS, Lambda, S3, SNS, SQS), Kubernetes, Docker, Terraform, Infrastructure as Code \\
\textbf{CI/CD \& Tools:} Azure DevOps, Git, GitHub, Visual Studio, CI/CD Pipelines, Automated Testing, Swagger/OpenAPI \\
\textbf{AI \& Automation:} AI Agents, Model Context Protocol (MCP), OpenAI API, Prompt Engineering, Agent Orchestration, REST/JSON \\
\textbf{Development \& Operations:} C\#, Python, .NET 8, ASP.NET Core, Microservices, Splunk, OpenTelemetry, Troubleshooting

\section*{Experience}
\resumeSubheading{Verisk Analytics}{Jan 2024 -- Present}{Software Engineer II}{Hyderabad, India}
\begin{itemize}
    \item Overhauled \textbf{Azure DevOps CI/CD} pipelines with staged approvals, parallel jobs, automated tests, and deployment gates, reducing release time by \textbf{50\%} and build cycles by \textbf{33\%}.
    \item Secured \textbf{ASP.NET Core} services on AWS EKS behind \textbf{Kong Gateway}, integrating \textbf{OAuth 2.0}, Okta-issued JWT validation, and authorization controls for protected cross-service communication.
    \item Applied \textbf{DevSecOps} practices across design, review, testing, and deployment, including input validation, least-privilege access, secure API contracts, and quality gates.
    \item Containerised workloads with \textbf{Docker}, orchestrated services on \textbf{Kubernetes/EKS}, and used \textbf{Terraform IaC} for repeatable environment provisioning and lifecycle management.
    \item Architected event-driven \textbf{AWS Lambda}, SNS, and SQS pipelines with idempotent processing, retry/backoff, and dead-letter handling, ensuring zero message loss during downstream failures.
    \item Built \textbf{LLM-powered workflow agents} with the \textbf{OpenAI API}, prompt engineering, and custom orchestration to automate Jira ticket creation and task routing across internal teams.
    \item Developed \textbf{MCP}-based tool integrations that exposed structured capabilities to agents, connecting model-driven workflows with external services through validated REST/JSON contracts.
    \item Co-designed \textbf{RulesRunner}, a stateless C\#/.NET microservice using CQRS and SQS that reduced client wait time from \textbf{3000ms to 500ms} (\textbf{83\%}) and enabled horizontal scaling.
    \item Instrumented services with \textbf{OpenTelemetry} and \textbf{Splunk}, accelerating troubleshooting of latency, authentication, deployment, and distributed pipeline failures.
\end{itemize}

\section*{Projects}
\resumeSubheading{Internal Workflow Automation Agent}{2025 -- 2026}{}{OpenAI API, MCP, C\#/Python, Jira, REST APIs}
\begin{itemize}
    \item Created an agent that interprets task requests and drafts Jira tickets using prompt-engineered OpenAI API calls, context management, validation, and custom workflow orchestration.
    \item Integrated MCP tools with downstream services and business rules, enabling reusable, controlled agent actions while reducing repetitive manual coordination.
\end{itemize}

\section*{Certifications}
\textbf{AWS Certified Cloud Practitioner} \hfill Amazon Web Services \;|\; Oct 2024

\section*{Education}
\textbf{B.Tech in Computer Engineering} --- Vishwakarma Institute of Technology, Pune (SPPU) \hfill 2024 \;|\; CGPA: 8.8

\end{document}
